\documentclass[12pt,a4paper]{article}

\usepackage{amsmath}
\usepackage{amssymb}
\usepackage{enumerate}
\usepackage{verbatim}
\usepackage[margin=1in]{geometry}
\usepackage{CJKutf8}

\begin{document}
\begin{CJK*}{UTF8}{bsmi}

\noindent \textbf{試題 :}

\vspace{0.5cm}

\section*{Problem 1 (15\%)}
Given an $n \times n$ matrix $A$, what is
\[
\|A\|_\infty = \max_{\|x\|_\infty = 1} \|Ax\|_\infty \ ?
\]
That is, write $\|A\|_\infty$ as an expression of $A$'s elements.

\vspace{0.5cm}

\section*{Problem 2 (25\%)}
Assume $A$ is symmetric and invertible. Formally prove that
\[
\text{cond}(A) \ge 1
\]
where $\text{cond}(\cdot)$ is the condition number. Assume 2-norm is used.

\vspace{0.25cm}
\noindent \textbf{Hint:} no need to check eigenvalues of matrices.

\vspace{0.5cm}

\section*{Problem 3 (20\%)}
Assume $L$ is a lower triangular square matrix stored in the compressed row format. That is, assume arrays \texttt{val}, \texttt{col\_ind}, and \texttt{row\_ptr} are already available. No need to worry about declaring variables and memory allocation. Write the MATLAB code to solve
\[
Lx = b
\]
We further make the following assumptions and requirements:
\begin{enumerate}[(1)]
    \item $L_{ii} \ne 0$, for all $i$.
    \item In \texttt{col\_ind}, indices of each row are in ascending order.
    \item Assume simple loops are used (It is like that you are writing Fortran/C/Java code though you use MATLAB syntax here and assume arrays start with index 1. Thus, you cannot use things such as \texttt{x(i:j)}, \texttt{find}, etc.).
    \item The complexity must be no more than $O(\text{nnz})$.
\end{enumerate}

\vspace{0.5cm}

\section*{Problem 4 (30\%)}
Assume a square matrix $A$ is stored in the compressed column format and we would like to find
\[
B = A^T .
\]
We store $B$ in compressed column format as well.

\vspace{0.25cm}
\noindent In the beginning, we try to find the number of non-zero elements in each column of $B$:

\begin{verbatim}
    for i=1:n+1
        bcol_ptr(i) = 0;
    nnz = acol_ptr(n+1)-1;
    for i=1:nnz
        bcol_ptr(arow_ind(i)+1) = bcol_ptr(arow_ind(i)+1) +1;
    bcol_ptr(1) = 1;
    for i=2:n+1
        bcol_ptr(i) = bcol_ptr(i) + bcol_ptr(i-1);
\end{verbatim}

\noindent Explain this segment of code and continue to finish the MATLAB code for obtaining $B$.

\vspace{0.25cm}
\noindent We further make the following assumptions and requirements:
\begin{enumerate}[(1)]
    \item Assume simple loops are used (It is like that you are writing Fortran/C/Java code though you use MATLAB syntax and assume arrays start with index 1. Thus, you cannot use things such as \texttt{x(i:j)}, \texttt{find}, etc.).
    \item The complexity must be no more than $O(\text{nnz})$.
    \item You can use only arrays \texttt{aval}, \texttt{arow\_ind}, \texttt{acol\_ptr}, \texttt{bval}, \texttt{bcol\_ind}, \texttt{brow\_ptr}. That is, you cannot allocate other arrays.
\end{enumerate}

\vspace{0.5cm}

\section*{Problem 5 (10\%)}
Given a linear system
\[
\begin{bmatrix}
1 & 2 & 1 \\
3 & 1 & -2 \\
1 & 2 & 1
\end{bmatrix}
x = 
\begin{bmatrix}
-1 \\ 1 \\ 2
\end{bmatrix}
\]
Use $x^{(0)} = \begin{bmatrix} 0 \\ 0 \\ 0 \end{bmatrix}$ as the initial point.

\begin{enumerate}[(a)]
    \item Do three iterations of the Jacobi method.
    \item Do three iterations of the Gauss-Seidel method.
\end{enumerate}

\end{CJK*}
\end{document}